\documentclass[aspectratio=169]{beamer}

% ==== Theme & basic look ====
\usetheme{Madrid}
\usecolortheme{seahorse}
\setbeamertemplate{navigation symbols}{}
\setbeamertemplate{footline}[frame number]


% Packages for mathematics  
\usepackage{amsmath}  
\usepackage{amsfonts}  
\usepackage{amssymb}  
\usepackage{CJKutf8}
% Support for \Tilde macro used throughout
\providecommand{\Tilde}[1]{\tilde{#1}}
\providecommand{\proofname}{Proof}
\newenvironment<>{proofs}[1][\proofname]{%
    \par
    \def\insertproofname{#1.}%
    \usebeamertemplate{proof begin}#2}
  {\usebeamertemplate{proof end}}
\makeatother

% Math shorthand and vector symbols
\newcommand{\grad}{\nabla}
\newcommand{\EE}{\mathbb{E}}
\newcommand{\dd}{\mathrm{d}}
\newcommand{\kbT}{k_B T}
\newcommand{\vx}{\mathbf{x}}

% Bold vector/time-indexed variables
\newcommand{\vxz}{\mathbf{x}_0}
\newcommand{\vxt}{\mathbf{x}_t}
\newcommand{\alpt}{\alpha_t}
\newcommand{\sigt}{\sigma_t}
\newcommand{\veps}{\boldsymbol{\epsilon}}
\newcommand{\vF}{\mathbf{F}}

% Title page details:   
\title[Large Foundation Models and Reasoning]{Large Foundation Models and Reasoning}  
\author{Pipi Hu}  
\institute[BIMSA]{Beijing Institute of Mathematical Sciences and Applications}  
\date{\today}  

\pgfdeclareimage[width=2cm]{logo}{../figures/logo.png}
\logo{\pgfuseimage{logo}}

% Outline slide (moved to later; avoid frame before \begin{document})

\setbeamertemplate{navigation symbols}{}
\AtBeginSection[]  
{  
    \begin{frame}<beamer>{Outline}         
        \tableofcontents[currentsection]  
    \end{frame}  
}  
  
\AtBeginSubsection[]  
{  
    \begin{frame}<beamer>{Outline}         
        \tableofcontents[currentsection,currentsubsection]  
    \end{frame}  
}  

% \setbeamertemplate{navigation symbols}{}
% \AtBeginSection[]
% {
%     \begin{frame}<beamer>{Outline}       \tableofcontents[currentsection]
%     \end{frame}
% }

\begin{document}  
  
% Title slide  
\begin{frame}  
  \titlepage  
\end{frame}  

\begin{frame}{Turing's words}
\centering
    The question of whether machines can think is like asking whether submarines can swim.
\end{frame}
  
% Presentation slides  

\section{Introduction to Reasoning in Large Foundation Models}

\begin{frame}{What is Reasoning?}
Reasoning is the process of drawing logical conclusions from given premises
\begin{itemize}
    \item \textbf{Deductive reasoning}: From general to specific
    \item \textbf{Inductive reasoning}: From specific to general  
    \item \textbf{Abductive reasoning}: Best explanation for observations
    \item \textbf{Common sense reasoning}: Everyday logical inference
\end{itemize}

\textcolor{red}{Key challenge}: How can large language models perform complex reasoning tasks?
\end{frame}

\begin{frame}{Types of Reasoning in AI}
\begin{enumerate}
    \item \textbf{Mathematical reasoning}: Solving equations, proofs
    \item \textbf{Logical reasoning}: Propositional and first-order logic
    \item \textbf{Causal reasoning}: Understanding cause-effect relationships
    \item \textbf{Commonsense reasoning}: Everyday knowledge and inference
    \item \textbf{Multi-step reasoning}: Chain of thought processes
\end{enumerate}

\textcolor{red}{Question}: Can foundation models truly reason or just pattern match?
\end{frame}

\section{Chain-of-Thought Reasoning}

\begin{frame}{Chain-of-Thought (CoT) Prompting}
Chain-of-Thought prompting encourages models to show their reasoning process
\begin{equation}
\text{Input} \rightarrow \text{Step 1} \rightarrow \text{Step 2} \rightarrow \cdots \rightarrow \text{Answer}
\end{equation}

\textbf{Example}:
\begin{itemize}
    \item Problem: "Sarah has 12 apples. She gives 3 to Tom and 4 to Mary. How many does she have left?"
    \item CoT: "Sarah starts with 12 apples. She gives 3 to Tom: 12 - 3 = 9. She gives 4 to Mary: 9 - 4 = 5. Sarah has 5 apples left."
\end{itemize}
\end{frame}

\begin{frame}{Mathematical Formulation of CoT}
Let $\mathcal{D} = \{(x_i, y_i)\}_{i=1}^N$ be a dataset where:
\begin{itemize}
    \item $x_i$: Input problem
    \item $y_i$: Ground truth answer
\end{itemize}

\textbf{CoT Training Objective}:
\begin{equation}
\mathcal{L}_{CoT} = \mathbb{E}_{(x,y) \sim \mathcal{D}} \left[ -\log p(y | x, \text{CoT}) \right]
\end{equation}

where CoT represents the intermediate reasoning steps.

\textbf{Key insight}: Decompose complex reasoning into simpler steps
\end{frame}

\begin{frame}{Benefits of Chain-of-Thought}
\begin{enumerate}
    \item \textbf{Transparency}: Makes reasoning process visible
    \item \textbf{Error localization}: Easier to identify where reasoning fails
    \item \textbf{Improved accuracy}: Especially for multi-step problems
    \item \textbf{Human interpretability}: Aligns with human problem-solving
\end{enumerate}

\textbf{Empirical evidence}: CoT improves performance on:
\begin{itemize}
    \item Mathematical word problems
    \item Logical reasoning tasks  
    \item Commonsense reasoning
    \item Symbolic manipulation
\end{itemize}
\end{frame}

\section{Mathematical Reasoning}

\begin{frame}{Mathematical Problem Solving}
Large language models can solve mathematical problems through:
\begin{enumerate}
    \item \textbf{Step-by-step decomposition}
    \item \textbf{Algebraic manipulation}
    \item \textbf{Logical inference}
    \item \textbf{Pattern recognition}
\end{enumerate}

\textbf{Example - Quadratic Equation}:
\begin{align}
x^2 + 5x + 6 &= 0 \\
(x + 2)(x + 3) &= 0 \\
x &= -2 \text{ or } x = -3
\end{align}

\textcolor{red}{Challenge}: Ensuring mathematical correctness in generated solutions
\end{frame}

\begin{frame}{Formal Mathematical Reasoning}
For formal mathematical proofs, we need:
\begin{equation}
\text{Given: } \Gamma \vdash \phi \text{ (premises)} \rightarrow \text{Prove: } \psi \text{ (conclusion)}
\end{equation}

\textbf{Proof strategies}:
\begin{itemize}
    \item Direct proof: $\Gamma \vdash \phi \rightarrow \psi$
    \item Proof by contradiction: $\Gamma, \neg\psi \vdash \bot$
    \item Proof by induction: Base case + inductive step
    \item Proof by cases: Exhaustive enumeration
\end{itemize}

\textbf{Key challenge}: Ensuring logical validity of each step
\end{frame}

\begin{frame}{Mathematical Reasoning Capabilities}
\textbf{What LLMs can do}:
\begin{itemize}
    \item Solve arithmetic problems
    \item Perform algebraic manipulations
    \item Understand mathematical concepts
    \item Generate mathematical explanations
\end{itemize}

\textbf{Current limitations}:
\begin{itemize}
    \item Inconsistent mathematical accuracy
    \item Difficulty with complex proofs
    \item Limited formal verification
    \item Hallucination of mathematical facts
\end{itemize}

\textcolor{red}{Research direction}: Improving mathematical reasoning reliability
\end{frame}

\section{Logical Reasoning}

\begin{frame}{Propositional Logic Reasoning}
Foundation models can perform logical inference using:
\begin{equation}
\text{Premises: } P_1, P_2, \ldots, P_n \vdash \text{Conclusion: } Q
\end{equation}

\textbf{Logical connectives}:
\begin{align}
\text{Negation: } &\neg P \\
\text{Conjunction: } &P \land Q \\
\text{Disjunction: } &P \lor Q \\
\text{Implication: } &P \rightarrow Q \\
\text{Biconditional: } &P \leftrightarrow Q
\end{align}

\textbf{Example}: If it rains, then the ground is wet. It's raining. Therefore, the ground is wet.
\end{frame}

\begin{frame}{First-Order Logic Reasoning}
For more complex reasoning involving quantifiers:
\begin{align}
\text{Universal: } &\forall x, P(x) \\
\text{Existential: } &\exists x, P(x)
\end{align}

\textbf{Logical inference rules}:
\begin{itemize}
    \item Modus Ponens: $P \rightarrow Q, P \vdash Q$
    \item Modus Tollens: $P \rightarrow Q, \neg Q \vdash \neg P$
    \item Universal Instantiation: $\forall x P(x) \vdash P(c)$
    \item Existential Generalization: $P(c) \vdash \exists x P(x)$
\end{itemize}

\textcolor{red}{Challenge}: Scaling logical reasoning to complex domains
\end{frame}

\begin{frame}{Commonsense Reasoning}
Everyday logical inference that humans perform naturally:
\begin{itemize}
    \item \textbf{Causal reasoning}: Understanding cause-effect relationships
    \item \textbf{Temporal reasoning}: Understanding time sequences
    \item \textbf{Spatial reasoning}: Understanding spatial relationships
    \item \textbf{Social reasoning}: Understanding social interactions
\end{itemize}

\textbf{Example}: "John went to the store. He bought milk." 
\begin{itemize}
    \item Implicit: John had money, store was open, milk was available
    \item Inference: John now has milk, spent money, was at store
\end{itemize}

\textcolor{red}{Key insight}: Commonsense requires world knowledge and inference
\end{frame}

\section{Advanced Reasoning Techniques}

\begin{frame}{Tree of Thoughts (ToT)}
Extends CoT by exploring multiple reasoning paths:
\begin{equation}
\text{Problem} \rightarrow \text{Multiple CoT paths} \rightarrow \text{Best solution}
\end{equation}

\textbf{Process}:
\begin{enumerate}
    \item Generate multiple reasoning paths
    \item Evaluate each path
    \item Select the best path
    \item Refine if necessary
\end{enumerate}

\textbf{Advantages}:
\begin{itemize}
    \item Explores multiple solution strategies
    \item Can backtrack from wrong paths
    \item More robust than single-path reasoning
\end{itemize}
\end{frame}

\begin{frame}{Self-Consistency}
Ensures reasoning consistency by:
\begin{equation}
\text{Multiple reasoning attempts} \rightarrow \text{Consensus answer}
\end{equation}

\textbf{Method}:
\begin{enumerate}
    \item Generate multiple reasoning paths for same problem
    \item Check consistency across paths
    \item Select most consistent answer
    \item Flag inconsistencies for review
\end{enumerate}

\textbf{Benefits}:
\begin{itemize}
    \item Reduces reasoning errors
    \item Identifies uncertain cases
    \item Improves overall reliability
\end{itemize}
\end{frame}

\begin{frame}{Program-Aided Language Models (PAL)}
Combines natural language reasoning with code execution:
\begin{equation}
\text{Problem} \rightarrow \text{Generate code} \rightarrow \text{Execute} \rightarrow \text{Answer}
\end{equation}

\textbf{Example}:
\begin{itemize}
    \item Problem: "Calculate the area of a circle with radius 5"
    \item Generated code: \texttt{import math; area = math.pi * 5**2}
    \item Execution: Returns 78.54
    \item Answer: "The area is approximately 78.54 square units"
\end{itemize}

\textbf{Advantages}:
\begin{itemize}
    \item Precise mathematical computation
    \item Verifiable results
    \item Handles complex calculations
\end{itemize}
\end{frame}

\section{Reasoning Evaluation}

\begin{frame}{Evaluating Reasoning Capabilities}
\textbf{Key metrics for reasoning evaluation}:
\begin{enumerate}
    \item \textbf{Accuracy}: Correctness of final answers
    \item \textbf{Consistency}: Logical coherence of reasoning steps
    \item \textbf{Completeness}: Coverage of all necessary steps
    \item \textbf{Efficiency}: Optimal reasoning path length
\end{enumerate}

\textbf{Evaluation datasets}:
\begin{itemize}
    \item GSM8K: Grade school math problems
    \item MATH: High school mathematics
    \item LogiQA: Logical reasoning
    \item CommonsenseQA: Commonsense reasoning
\end{itemize}
\end{frame}

\begin{frame}{Reasoning Error Analysis}
\textbf{Common reasoning errors}:
\begin{enumerate}
    \item \textbf{Logical fallacies}: Invalid inference patterns
    \item \textbf{Factual errors}: Incorrect knowledge application
    \item \textbf{Calculation mistakes}: Arithmetic errors
    \item \textbf{Incomplete reasoning}: Missing steps
\end{enumerate}

\textbf{Error detection methods}:
\begin{itemize}
    \item Consistency checking across multiple attempts
    \item Step-by-step verification
    \item External knowledge validation
    \item Human expert review
\end{itemize}
\end{frame}

\section{Future Directions}

\begin{frame}{Challenges in Reasoning}
\textbf{Current limitations}:
\begin{enumerate}
    \item \textbf{Hallucination}: Generating plausible but incorrect reasoning
    \item \textbf{Inconsistency}: Contradictory reasoning within same response
    \item \textbf{Context limitations}: Difficulty with very long reasoning chains
    \item \textbf{Domain specificity}: Poor generalization across domains
\end{enumerate}

\textbf{Research directions}:
\begin{itemize}
    \item Improving reasoning reliability
    \item Scaling to more complex problems
    \item Better evaluation methodologies
    \item Integration with external tools
\end{itemize}
\end{frame}

\begin{frame}{Emerging Approaches}
\textbf{Promising research directions}:
\begin{enumerate}
    \item \textbf{Neuro-symbolic reasoning}: Combining neural and symbolic methods
    \item \textbf{Multi-agent reasoning}: Collaborative problem solving
    \item \textbf{Interactive reasoning}: Human-AI reasoning collaboration
    \item \textbf{Formal verification}: Mathematical proof of reasoning correctness
\end{enumerate}

\textbf{Key questions}:
\begin{itemize}
    \item Can we achieve truly reliable reasoning?
    \item How to scale reasoning to complex domains?
    \item What is the role of human oversight?
\end{itemize}
\end{frame}

\section{Conclusion}

\begin{frame}{Summary}
\textbf{Key insights about reasoning in foundation models}:
\begin{enumerate}
    \item Chain-of-thought prompting significantly improves reasoning
    \item Mathematical and logical reasoning show promising results
    \item Commonsense reasoning remains challenging
    \item Evaluation and error detection are crucial
\end{enumerate}

\textbf{Future outlook}:
\begin{itemize}
    \item Continued improvement in reasoning capabilities
    \item Better integration with external tools
    \item More robust evaluation methods
    \item Potential for truly reliable AI reasoning
\end{itemize}
\end{frame}

\begin{frame}{Open Questions}
\textbf{Research questions for the community}:
\begin{enumerate}
    \item How can we ensure reasoning reliability at scale?
    \item What are the fundamental limits of neural reasoning?
    \item How to best combine symbolic and neural approaches?
    \item What role should human oversight play?
\end{enumerate}

\textbf{Contact}: For questions and collaboration, please reach out at {\color{red}(\href{mailto:hpp@bimsa.cn}{hpp@bimsa.cn})}
\end{frame}

% Thank you slide  
\begin{frame}{Thank you}  
  \centering  
  Questions and Discussion
  \\[2em]
  Thank you for your attention!  
\end{frame}  
  
\end{document}
